\subsection{Сила Архимеда}
%%%%%%%%%%%%%%%%%%%%%%%%%%%%%%%%%%%%%%%%%%%%%%%%%%%%
%%%%%%%%%%%%%%%%%%%%%%%%%%%%%%%%%%%%%%%%%%%%%%%%%%%%
\z{В сплошном кубике из дерева масса $100 \,\text{г}$, объём $120 \,\text{см}^3$. Какой массы кусок алюминия нужно к нему прикрепить, чтобы он утонул?
}
%
{$32 \,\text{г}$}

%%%%%%%%%%%%%%%%%%%%%%%%%%%%%%%%%%%%%%%%%%%%%%%%%%%%
\z{На очень точных весах гирька из алюминия массой $50 \,\text{г}$ уравновешена гирьками из неизвестного металла, общая масса которых $49.99 \,\text{г}$. Найти плотность неизвестного металла.
}
%
{$\rho_x = \frac{\rho_{\text{возд}} m_2}{m_2 - m_1 + \frac{\rho_{\text{возд}} m_1}{\rho_{\text{Al}}}} \approx 4600 \,\text{кг/м}^3$}

%%%%%%%%%%%%%%%%%%%%%%%%%%%%%%%%%%%%%%%%%%%%%%%%%%%%
\z{Гладкий кубик со стороной $a$ лежит на дне гладкого сосуда с водой глубиной $h$. Чтобы оторвать кубик от дна, потребовалась сила $F$. Чему равна плотность кубика?
}
%
{$\rho = \frac{F - \rho_{\text{в}} g (h-a) a^2}{g a^3}$}


%%%%%%%%%%%%%%%%%%%%%%%%%%%%%%%%%%%%%%%%%%%%%%%%%%%%
\z[{\begin{wrapfigure}[3]{r}{0.25\textwidth}%
		\vspace{-20pt}%
		\begin{flushright}%
 			\includegraphics[width=0.25\textwidth]{4.3.4}%
		\end{flushright}%
	\end{wrapfigure}%
	}]%
{На границе раздела двух жидкостей плотностей $\rho_1$ и $\rho_2$ плавает шайба плотности $\rho$ ($\rho_1 < \rho < \rho_2$). Высота шайбы $H$. Определите глубину её погружения во вторую жидкость.
\vspace{2em}
}
%
{$h = H \frac{\rho - \rho_1}{\rho_2 - \rho_1}$}

%%%%%%%%%%%%%%%%%%%%%%%%%%%%%%%%%%%%%%%%%%%%%%%%%%%%
\z[{\begin{wrapfigure}[3]{r}{0.15\textwidth}%
		\vspace{-20pt}%
		\begin{flushright}%
 			\includegraphics[width=0.15\textwidth]{4.3.5}%
		\end{flushright}%
	\end{wrapfigure}%
	}]%
{Однородный поплавок удерживается в воде при помощи закрепленной на дне стакана нити. Найдите плотность поплавка, если известно, что сила натяжения нити в 4 раза меньше силы тяжести, действующей на шарик. 
\vspace{2em}
}
%
{$\rho_\text{попл}=\frac{4}{5}\rho_\text{в}= 800 \, \text{кг/м}^3$}

%%%%%%%%%%%%%%%%%%%%%%%%%%%%%%%%%%%%%%%%%%%%%%%%%%%%
\z[{\begin{wrapfigure}[4]{r}{0.15\textwidth}%
		\vspace{-20pt}%
		\begin{flushright}%
 			\includegraphics[width=0.15\textwidth]{4.3.6}%
		\end{flushright}%
	\end{wrapfigure}%
	}]%
{Определите силу натяжения нити, связывающей два шарика объёмом $10 \,\text{см}^3$ каждый, если верхний шарик плавает, наполовину погрузившись в воду, а нижний шарик в три раза тяжелее верхнего.
\vspace{2em}
}
%
{$F = \frac{1}{8} \rho g V = 1.2 \cdot 10^{-2} \,\text{Н}$}

%%%%%%%%%%%%%%%%%%%%%%%%%%%%%%%%%%%%%%%%%%%%%%%%%%%%
\z[{\begin{wrapfigure}[5]{r}{0.15\textwidth}%
		\vspace{-20pt}%
		\begin{flushright}%
 			\includegraphics[width=0.15\textwidth]{4.3.7}%
		\end{flushright}%
	\end{wrapfigure}%
	}]%
{В цилиндрическом сосуде с водой плавает поплавок, к которому привязан груз массой $3 \,\text{кг}$ и объёмом $1 \,\text{л}$. Как и насколько изменится уровень воды в сосуде, если нить оборвётся и груз упадёт на дно? Площадь дна сосуда $50 \,\text{см}^2$.
\vspace{2em}
}
%
{$\Delta h = \frac{m - \rho V}{\rho S} \quad\Rightarrow\quad -40 \,\text{см}$}
